\section{AI 对软件工程的影响：更需要 Systems Judgment}

课程最后转向教育和 IDE 未来。Asif 的判断不是“CS 基础失效”，而是生成代码增多后，理解 state、performance、security、failure 和 interface 的能力更重要。Agent 可以写更多代码，也会更快制造依赖、重试和权限问题；工程师的工作从逐行输入扩展到定义 contract、验证 evidence 和管理 change。

\subsection{IDE 变成 Agent workspace}

IDE 会承载多个 agent、background task、terminal、browser、test 和 deployment feedback。基础设施必须让用户知道 agent 在读什么、改什么、运行什么，并保留可暂停、可审查、可回滚的 control plane。更强 autonomy 不应以失去 workspace state visibility 为代价。

\begin{importantbox}{AI-native software engineering 的四项核心能力}
\begin{enumerate}
  \item 设计 versioned interface、invariant 和 permission；
  \item 为 generated change 建立 test、eval、diff 和 rollback；
  \item 识别 queue、database、network 与 provider 的系统瓶颈；
  \item 把 incident 和 user feedback 转成 guardrail，而不是只修 prompt。
\end{enumerate}
\end{importantbox}

\subsection{AI-assisted incident response}

AI 可以汇总 logs、关联 deploy、生成 query、解释 runbook 和提出 hypothesis，但生产事故处于证据不完整、成本快速变化的环境。AI 不应直接拥有无限 production authority；incident commander 需要确认目标、停止 feedback loop、选择风险和记录决策。

\lecturefigure{16-incident-learning.png}{AI 可以压缩证据处理，但 human incident command 仍负责生产决策。}{本地字幕 47:50--48:37；概念重绘。}

\begin{knowledgebox}{读图：Assistant 最适合缩短哪一步}
它可以从海量 metrics/logs 中定位时间相关性，生成 safe read-only query，比较最近 deploy，或把聊天记录整理成 timeline。高风险 write、failover、data repair 和 customer communication 仍需要显式 owner 与 approval。
\end{knowledgebox}

\begin{warningbox}{AI 生成的根因同样需要证据}
语言模型容易把相关性写成因果故事。每个 hypothesis 都应链接 metric、trace、log、config 或 reproduction；postmortem 记录 unknown 与 rejected hypothesis，避免把流畅解释误当事实。
\end{warningbox}

\teachervoice{课堂的结尾把 AI-assisted incident response 当成可能方向。最稳妥的迁移结论是：让 AI 提高 evidence bandwidth，让人保留目标、风险和 authority。}

\subsection{本章小结}

AI coding 并未降低 systems knowledge 的价值。相反，生成能力越强，验证、隔离、恢复、安全和资源会计越成为产品与职业的核心。

